\documentclass[10pt]{article}
\usepackage[a4paper,margin=25mm]{geometry}
\usepackage{amsmath,amssymb,amsthm,booktabs}
\newtheorem{proposition}{Proposition}
\title{The regular columns of the small-group atlas}
\author{Human Compression project}\date{}
\begin{document}\maketitle
We consider the 79 group models printed in Tables 2--7 of Hartley,
Hubard and Leemans, \emph{Two atlases of abstract chiral polytopes for
small groups}, Ars Mathematica Contemporanea 5 (2012), 371--382.
There are 88 table occurrences, since some groups occur more than once.
A regular marking of rank $r$ means an ordered tuple of nonidentity
involutions $\rho_0,\ldots,\rho_{r-1}$ which generates the group, satisfies
$\rho_i\rho_j=\rho_j\rho_i$ for $|i-j|>1$, and has the intersection property
\[
 \langle\rho_i:i\in I\rangle\cap\langle\rho_j:j\in J\rangle
 =\langle\rho_k:k\in I\cap J\rangle
 \qquad(I,J\subseteq\{0,\ldots,r-1\}).
\]
Markings are identified by automorphisms carrying each generator to the
corresponding generator. Reversal is not additionally identified. The
standard correspondence with regular abstract polytopes is retained;
we do not assert a classification under a different incidence convention.

\begin{proposition}
Relative to the acquired literal models, group-operation implementations,
and the reviewed enumeration theorem, the regular columns for all 79
printed groups are determined by the accompanying census files. Of the
88 printed rank-count rows, 87 agree with these results. In Table 4,
page 379, the rank-three entry for $\mathrm{PSL}_2(27)$ is 26, whereas
the source prints 27. All other regular entries agree rank by rank.
\end{proposition}
\begin{proof}
Fourteen models are supplied by the already reviewed complete regular
censuses; their rank lists are acquired without repeating their searches.
For the remaining 65 models, the acquired enumeration constructs
involutory prefixes up to the complete automorphism action. At an
extension, representatives of the relevant double cosets enumerate all
inequivalent choices in the required centralizers. The intersection test
and the acquired normal-closure rejection remove only impossible prefixes.
A string C-group tuple is independent: removing a generator strictly
reduces the subgroup it generates. The increasing subgroup chain therefore
bounds its rank by $\lfloor\log_2|G|\rfloor$. These are the hypotheses and
conclusion of the previously reviewed enumeration theorem; this application
uses that theorem without altering its search or rejection rules.

The action on each literal model is acquired from the reviewed action
construction. For a normalizing permutation group $A$, the equality
$|A:C_A(G)|=|\operatorname{Aut}G|$ proves that its conjugation action is
complete. Where necessary, the construction instead uses the faithful
common action of $G$ and its automorphism group on the involutions.
The checked generator map connects that action to the original model.
Both constructions and their faithfulness hypotheses are retained inputs.

Every one of the 65 actual executions returned its complete rank list,
with a clean terminal result and unchanged input files. The independent
receipt verifier reconstructs the complete input, driver and output
inventories and reparses the returned census. It neither executes another
search nor reads the source counts as search inputs. Joining these 65
lists with the 14 retained lists gives exactly the 79 printed group keys,
without overlap. Joining them to all 88 source occurrences gives the
stated comparison. The machine-readable result records every rank list,
including the source discrepancy and all duplicate joins.
\end{proof}

\paragraph{An independent rank-three check for $\mathrm{PSL}_2(27)$.}
This count also has a short finite verification which does not use the
prefix enumeration. Write $G=\mathrm{PSL}_2(27)$ in its natural action on
28 points. Its order is 9828. The accompanying script checks literal
equality with the standard constructor, and acquires a normalizing action
of order 58968 inducing all automorphisms. There are 351 involutions in
$G$, forming one orbit. Fix an involution $x$. The script lists all 14
distinct involutions commuting with $x$ and verifies that the stabilizer
of $x$ is transitive on them. Thus all possible ordered commuting endpoint
pairs $(x,z)$ are equivalent. The stabilizer of one such pair has order 12.

For this fixed pair, the middle generator $y$ must be one of the 351
involutions. Exactly 312 choices generate $G$ together with $x,z$.
For every one of them the explicit subgroup computation gives
\[
 \langle x,y\rangle\cap\langle y,z\rangle=\langle y\rangle.
\]
This is the rank-three intersection criterion: the rank-two subgroups
are dihedral string C-groups, and their intersection is the common
rank-one subgroup. Equivalently, the other rank-two intersections follow
from $\langle x,z\rangle\cong C_2^2$, the displayed equality, and generation.
The endpoint stabilizer acts freely on the 312 possibilities, because an
automorphism fixing the three generating elements is the identity.
Consequently the number of rank-three markings is $312/12=26$.
The script separately checks its orbits and their sizes. Enumeration of
all elements, subgroup operations and the automorphism-group order are
explicit GAP premises of this independent finite check. This paragraph
makes no assertion about higher ranks on its own.

\paragraph{Scope and attribution.}
The source determines the group roster and the marking convention.
The previous complete regular censuses, the enumeration theorem and the
common-action construction supply the reused mathematics. The new work
here is their complete application to the 65 remaining models, the join
to all printed regular entries, and the separate fixed-endpoint check of
the exceptional row. This is a correction of a finite published count,
not a claimed solution of an open conjecture or a claim of priority.
The chiral columns, first atlas, and the scope of the historical ATLAS
roster require their separate arguments. No whole-manuscript completion,
Lean verification or compression measurement follows from this component.
\end{document}
